%==============================================================================
% Chapitre 20 - Optiques et viseurs modernes
%==============================================================================

\section{Introduction}

Les systèmes de visée optiques permettent d'améliorer la capacité du tireur à
détecter, identifier et engager une cible.

Ils complètent ou remplacent les organes de visée mécaniques étudiés au
chapitre précédent.

Leur développement a profondément transformé les performances des systèmes
d'armes modernes.

\begin{aretenir}

Les systèmes optiques n'améliorent pas directement la précision intrinsèque
d'une arme, mais ils permettent au tireur d'exploiter plus efficacement cette
précision.

\end{aretenir}

\section{Principes généraux}

Un système optique a pour objectif de faciliter l'acquisition et le pointage
de la cible.

Selon les technologies utilisées, il peut :

\begin{itemize}
    \item grossir l'image ;
    \item afficher un réticule ;
    \item mesurer une distance ;
    \item compenser certains paramètres balistiques ;
    \item permettre l'observation de nuit.
\end{itemize}

\section{La ligne de visée}

Comme pour les organes mécaniques, l'axe optique est généralement situé au
dessus de l'axe du canon.

Cette différence de hauteur doit être prise en compte lors du réglage de
l'arme.

Elle explique notamment l'existence :

\begin{itemize}
    \item du point sécant proche ;
    \item du point sécant lointain ;
    \item des écarts observés aux très courtes distances.
\end{itemize}

\section{Lunettes de tir}

La lunette de tir est l'un des systèmes optiques les plus répandus.

Elle permet :

\begin{itemize}
    \item d'agrandir l'image de la cible ;
    \item d'améliorer la précision du pointage ;
    \item d'effectuer certaines corrections balistiques.
\end{itemize}

Une lunette est généralement constituée :

\begin{itemize}
    \item d'un objectif ;
    \item d'un système de lentilles ;
    \item d'un réticule ;
    \item d'un oculaire ;
    \item de mécanismes de réglage.
\end{itemize}

\section{Grossissement}

Le grossissement indique combien de fois l'image est agrandie.

\begin{exemple}

Une lunette de grossissement :

\begin{equation}
10\times
\end{equation}

présente une image apparaissant dix fois plus proche.

\end{exemple}

Les lunettes modernes peuvent être :

\begin{itemize}
    \item à grossissement fixe ;
    \item à grossissement variable.
\end{itemize}

\section{Champ de vision}

Le champ de vision représente la portion de terrain visible à travers
l'optique.

En règle générale :

\begin{itemize}
    \item un fort grossissement réduit le champ de vision ;
    \item un faible grossissement augmente le champ de vision.
\end{itemize}

Le choix dépend donc du compromis recherché.

\section{Le réticule}

Le réticule constitue la référence visuelle utilisée pour viser.

Il peut prendre de nombreuses formes :

\begin{itemize}
    \item croix simple ;
    \item duplex ;
    \item mil-dot ;
    \item grille graduée ;
    \item réticules balistiques spécialisés.
\end{itemize}

Le réticule peut être placé :

\begin{itemize}
    \item dans le premier plan focal ;
    \item dans le second plan focal.
\end{itemize}

\section{Premier et second plan focal}

Dans une lunette à premier plan focal (FFP) :

\begin{itemize}
    \item le réticule change de taille avec le grossissement ;
    \item les graduations restent valides à tous les grossissements.
\end{itemize}

Dans une lunette à second plan focal (SFP) :

\begin{itemize}
    \item le réticule conserve une taille apparente constante ;
    \item certaines graduations ne sont exactes qu'à un grossissement donné.
\end{itemize}

\section{Réglages en site et en dérive}

Les lunettes permettent généralement deux réglages principaux :

\begin{itemize}
    \item le site (élévation) ;
    \item la dérive.
\end{itemize}

Ces réglages sont effectués à l'aide de tourelles graduées.

Ils permettent de compenser :

\begin{itemize}
    \item la chute du projectile ;
    \item le vent ;
    \item certains écarts systématiques.
\end{itemize}

\section{MOA et milliradians}

Les corrections angulaires sont généralement exprimées :

\begin{itemize}
    \item en MOA (Minute Of Angle) ;
    \item en milliradians (mil).
\end{itemize}

Ces unités seront étudiées plus en détail dans un chapitre ultérieur consacré
aux mesures angulaires.

\begin{aretenir}

Les lunettes corrigent des angles et non directement des distances.

\end{aretenir}

\section{Parallaxe}

La parallaxe apparaît lorsque le plan du réticule et celui de l'image ne
coïncident pas parfaitement.

Elle peut provoquer un déplacement apparent du point visé lorsque l'œil se
déplace derrière l'oculaire.

Les lunettes modernes disposent souvent d'un réglage de parallaxe.

\section{Viseurs à point rouge}

Les viseurs à point rouge projettent un repère lumineux permettant une
acquisition rapide de la cible.

Ils présentent plusieurs avantages :

\begin{itemize}
    \item rapidité ;
    \item simplicité ;
    \item large champ de vision ;
    \item utilisation intuitive.
\end{itemize}

Ils sont particulièrement adaptés aux engagements à courte et moyenne
distance.

\section{Collimateurs}

Les collimateurs sont conçus pour présenter un repère optique apparaissant à
une distance virtuelle très importante.

Le tireur peut ainsi observer simultanément :

\begin{itemize}
    \item le réticule ;
    \item la cible.
\end{itemize}

Cette technologie est largement utilisée sur les armes individuelles et les
véhicules.

\section{Télémètres}

Certains systèmes optiques intègrent un télémètre.

Le plus souvent, il s'agit d'un télémètre laser.

Sa fonction est de mesurer la distance jusqu'à la cible.

Cette information peut ensuite être utilisée pour :

\begin{itemize}
    \item calculer une correction ;
    \item alimenter un calculateur balistique ;
    \item assister le tireur.
\end{itemize}

\section{Vision nocturne}

Les dispositifs de vision nocturne permettent l'observation dans des
conditions de faible luminosité.

On distingue notamment :

\begin{itemize}
    \item l'intensification de lumière ;
    \item l'imagerie thermique.
\end{itemize}

Ces technologies reposent sur des principes physiques différents.

\section{Imagerie thermique}

Les systèmes thermiques détectent le rayonnement infrarouge émis par les
objets.

Ils permettent notamment :

\begin{itemize}
    \item la détection ;
    \item l'observation ;
    \item l'identification dans certaines conditions.
\end{itemize}

Ils sont largement utilisés dans les systèmes militaires modernes.

\section{Optiques et simulation}

Les simulateurs modernes reproduisent souvent :

\begin{itemize}
    \item les grossissements ;
    \item les réticules ;
    \item les limitations optiques ;
    \item les effets de parallaxe ;
    \item les dispositifs thermiques.
\end{itemize}

Cette reproduction contribue fortement à l'immersion et au réalisme.

\begin{simulation}

Dans les systèmes d'entraînement avancés, l'optique constitue parfois un
modèle indépendant relié au moteur balistique et au moteur graphique.

\end{simulation}

\section{Vers les systèmes de conduite de tir}

Les systèmes optiques modernes sont de plus en plus associés à des aides
électroniques :

\begin{itemize}
    \item calculateurs balistiques ;
    \item télémètres ;
    \item capteurs météorologiques ;
    \item systèmes de conduite de tir.
\end{itemize}

Ces systèmes seront étudiés dans les prochains chapitres.

\section{À retenir}

\begin{aretenir}

\begin{itemize}
    \item Les optiques améliorent la capacité du tireur à viser.
    \item Le grossissement et le champ de vision sont liés.
    \item Les réticules peuvent intégrer des informations balistiques.
    \item Les corrections sont généralement exprimées en MOA ou en mil.
    \item Les systèmes modernes associent souvent optique et électronique.
\end{itemize}

\end{aretenir}
